% Template LaTeX pour le Rapport TP1 - Réservation de Terrains
\documentclass[12pt,a4paper]{article}

% Packages
\usepackage[utf8]{inputenc}
\usepackage[french]{babel}
\usepackage{geometry}
\usepackage{graphicx}
\usepackage{listings}
\usepackage{xcolor}
\usepackage{hyperref}
\usepackage{fancyhdr}
\usepackage{titlesec}
\usepackage{enumitem}

% Configuration
\geometry{margin=2.5cm}
\hypersetup{
    colorlinks=true,
    linkcolor=blue,
    urlcolor=blue,
    citecolor=blue
}

% Style du code
\definecolor{codegreen}{rgb}{0,0.6,0}
\definecolor{codegray}{rgb}{0.5,0.5,0.5}
\definecolor{codepurple}{rgb}{0.58,0,0.82}
\definecolor{backcolour}{rgb}{0.95,0.95,0.92}

\lstdefinestyle{mystyle}{
    backgroundcolor=\color{backcolour},   
    commentstyle=\color{codegreen},
    keywordstyle=\color{magenta},
    numberstyle=\tiny\color{codegray},
    stringstyle=\color{codepurple},
    basicstyle=\ttfamily\footnotesize,
    breakatwhitespace=false,         
    breaklines=true,                 
    captionpos=b,                    
    keepspaces=true,                 
    numbers=left,                    
    numbersep=5pt,                  
    showspaces=false,                
    showstringspaces=false,
    showtabs=false,                  
    tabsize=2
}
\lstset{style=mystyle}

% En-tête et pied de page
\pagestyle{fancy}
\fancyhf{}
\lhead{TP1 - Système de Réservation}
\rhead{\today}
\cfoot{\thepage}

% Début du document
\begin{document}

% Page de garde
\begin{titlepage}
    \centering
    \vspace*{2cm}
    
    {\Huge\bfseries Système de Réservation\\de Terrains de Football\par}
    \vspace{1cm}
    {\Large Travail Pratique 1\par}
    \vspace{2cm}
    
    {\Large\itshape Votre Nom\par}
    {\large Matricule: XXXXXXXX\par}
    \vspace{1cm}
    
    {\large Cours: [Nom du cours]\par}
    {\large Professeur: [Nom du professeur]\par}
    \vspace{2cm}
    
    {\large Université du Québec à Rimouski (UQAR)\par}
    \vspace{1cm}
    
    {\large\today\par}
    
    \vfill
\end{titlepage}

% Table des matières
\tableofcontents
\newpage

% ================================
% 1. INTRODUCTION
% ================================
\section{Introduction}

\subsection{Contexte du projet}
Ce projet consiste en la réalisation d'un système de réservation de terrains de football permettant aux clients de réserver des créneaux horaires, d'effectuer des paiements sécurisés et de télécharger leurs factures.

\subsection{Objectifs}
Les objectifs principaux de ce travail sont :
\begin{itemize}
    \item Développer une application web MVC complète avec ASP.NET Core
    \item Implémenter un système d'authentification multi-rôles
    \item Intégrer un système de paiement électronique (Stripe)
    \item Gérer la persistance des données avec Entity Framework Core
    \item Créer des interfaces utilisateur modernes et responsives
\end{itemize}

\subsection{Technologies utilisées}
\begin{itemize}
    \item \textbf{Backend:} ASP.NET Core MVC 8.0, C\#
    \item \textbf{Base de données:} SQL Server (LocalDB), Entity Framework Core 9.0
    \item \textbf{Authentification:} ASP.NET Core Identity
    \item \textbf{Paiement:} Stripe SDK .NET
    \item \textbf{Frontend:} Bootstrap 5, Font Awesome, jQuery
    \item \textbf{Outils:} Visual Studio/VS Code, Git, Postman
\end{itemize}

% ================================
% 2. ARCHITECTURE
% ================================
\section{Architecture du système}

\subsection{Architecture MVC}
L'application suit le pattern Model-View-Controller (MVC) :
\begin{itemize}
    \item \textbf{Models:} 7 entités principales + ViewModels
    \item \textbf{Views:} 25+ vues Razor avec Bootstrap 5
    \item \textbf{Controllers:} 6 contrôleurs principaux
\end{itemize}

\subsection{Modèle de données}

\subsubsection{Entités principales}
\begin{enumerate}
    \item \textbf{Utilisateur} - Gestion des comptes (Client, Fournisseur, Admin)
    \item \textbf{Terrain} - Informations sur les terrains de football
    \item \textbf{Creneau} - Créneaux horaires disponibles
    \item \textbf{PanierItem} - Articles dans le panier
    \item \textbf{Reservation} - Réservations effectuées
    \item \textbf{Paiement} - Transactions Stripe
    \item \textbf{Facture} - Factures générées
\end{enumerate}

\subsubsection{Relations entre entités}
% Insérer un diagramme ER ici
\begin{center}
\textit{[Insérer le diagramme de la base de données]}
\end{center}

\subsection{Architecture des services}
Trois services métier principaux :
\begin{itemize}
    \item \textbf{CreneauService:} Gestion des créneaux et disponibilités
    \item \textbf{PaiementService:} Intégration avec Stripe
    \item \textbf{FactureService:} Génération et gestion des factures
\end{itemize}

% ================================
% 3. FONCTIONNALITÉS
% ================================
\section{Fonctionnalités implémentées}

\subsection{Authentification et autorisation}
\subsubsection{Gestion des utilisateurs}
\begin{itemize}
    \item Inscription avec validation des données
    \item Connexion avec option "Se souvenir de moi"
    \item Gestion du profil (modification nom, email, mot de passe)
    \item Déconnexion sécurisée
\end{itemize}

\subsubsection{Système de rôles}
Trois rôles distincts avec permissions spécifiques :
\begin{enumerate}
    \item \textbf{Client:} Parcourir, réserver, payer, consulter factures
    \item \textbf{Fournisseur:} Accès aux statistiques de revenus
    \item \textbf{Admin:} Dashboard complet, gestion des réservations/terrains/utilisateurs
\end{enumerate}

\subsection{Parcours client}

\subsubsection{Navigation et filtrage}
\begin{itemize}
    \item Affichage paginé des créneaux disponibles
    \item Filtres multiples :
    \begin{itemize}
        \item Par date
        \item Par type de terrain (5-a-side, 7-a-side, 11-a-side)
        \item Par localisation
    \end{itemize}
    \item Barre de recherche par nom de terrain
    \item Combinaison de plusieurs filtres simultanément
\end{itemize}

\subsubsection{Détails des créneaux}
Page détaillée affichant :
\begin{itemize}
    \item Image du terrain
    \item Informations complètes (nom, type, localisation)
    \item Horaires et disponibilités
    \item Prix par place
    \item Description du terrain
\end{itemize}

\subsubsection{Gestion du panier}
\begin{itemize}
    \item Ajout de créneaux multiples
    \item Modification des quantités
    \item Suppression d'articles
    \item Calcul automatique du total
    \item Validation des places disponibles
    \item Persistance du panier par utilisateur
\end{itemize}

\subsection{Paiement électronique}

\subsubsection{Intégration Stripe}
\begin{itemize}
    \item Création de PaymentIntent côté serveur
    \item Interface de paiement sécurisée (Stripe Elements)
    \item Gestion des erreurs de paiement
    \item Confirmation du paiement
    \item Traçabilité dans le dashboard Stripe
\end{itemize}

\subsubsection{Sécurité des transactions}
\begin{itemize}
    \item Clés API en variables d'environnement
    \item Mode sandbox pour les tests
    \item Validation côté serveur
    \item Protection CSRF avec antiforgery tokens
\end{itemize}

\subsection{Système de facturation}
\begin{itemize}
    \item Génération automatique après paiement réussi
    \item Numéro de facture unique (format: FAC-YYYYMMDD-XXXXXX)
    \item Consultation et téléchargement des factures
    \item Historique complet des factures
    \item Fonction d'impression optimisée
\end{itemize}

\subsection{Dashboard administrateur}

\subsubsection{Indicateurs clés (KPIs)}
\begin{itemize}
    \item Total des réservations
    \item Réservations confirmées
    \item Revenu total généré
    \item Taux de conversion
\end{itemize}

\subsubsection{Statistiques avancées}
\begin{itemize}
    \item Top 10 des clients (par dépenses)
    \item Taux de remplissage par terrain
    \item Réservations récentes
    \item Revenus par jour (7 derniers jours)
\end{itemize}

\subsubsection{Gestion}
\begin{itemize}
    \item Vue d'ensemble des réservations
    \item Gestion des terrains
    \item Gestion des utilisateurs
\end{itemize}

% ================================
% 4. IMPLÉMENTATION
% ================================
\section{Détails d'implémentation}

\subsection{Configuration de la base de données}
\subsubsection{Connection string}
\begin{lstlisting}[language=json]
"ConnectionStrings": {
  "DefaultConnection": "Server=(localdb)\\mssqllocaldb;
  Database=ReservationTerrainsDB;
  Trusted_Connection=True;
  MultipleActiveResultSets=true"
}
\end{lstlisting}

\subsubsection{Migrations}
Commandes utilisées :
\begin{lstlisting}[language=bash]
dotnet ef migrations add InitialCreate
dotnet ef database update
\end{lstlisting}

\subsection{Seeding de données}
Le fichier \texttt{DbInitializer.cs} peuple automatiquement la base avec :
\begin{itemize}
    \item 3 rôles (Admin, Client, Fournisseur)
    \item 3 utilisateurs de test
    \item 3 terrains différents
    \item 126 créneaux (3 terrains × 6 créneaux/jour × 7 jours)
\end{itemize}

\subsection{Extraits de code significatifs}

\subsubsection{Service de paiement}
\begin{lstlisting}[language={[Sharp]C}]
public async Task<string> CreerPaymentIntentAsync(
    decimal montant, string devise = "eur")
{
    var options = new PaymentIntentCreateOptions
    {
        Amount = (long)(montant * 100),
        Currency = devise,
        AutomaticPaymentMethods = new 
            PaymentIntentAutomaticPaymentMethodsOptions
        {
            Enabled = true,
        },
    };
    
    var service = new PaymentIntentService();
    var paymentIntent = await service.CreateAsync(options);
    return paymentIntent.ClientSecret;
}
\end{lstlisting}

% [Ajouter d'autres extraits pertinents]

% ================================
% 5. TESTS
% ================================
\section{Tests et validation}

\subsection{Tests fonctionnels}
\subsubsection{Scénarios testés}
\begin{enumerate}
    \item Navigation et filtrage des créneaux
    \item Inscription et connexion
    \item Ajout au panier et modification
    \item Processus de paiement complet
    \item Génération de factures
    \item Accès aux dashboards (Client, Fournisseur, Admin)
    \item Protection des routes par rôle
\end{enumerate}

\subsection{Tests de sécurité}
\begin{itemize}
    \item Validation des formulaires (client et serveur)
    \item Protection CSRF
    \item Autorisation par rôle
    \item Sanitisation des entrées
\end{itemize}

\subsection{Tests avec Postman}
% [Décrire les tests API effectués]

\subsection{Résultats des tests}
Tableau récapitulatif :

\begin{center}
\begin{tabular}{|l|c|}
\hline
\textbf{Fonctionnalité} & \textbf{Statut} \\
\hline
Authentification & ✓ Réussi \\
Filtrage créneaux & ✓ Réussi \\
Panier multi-créneaux & ✓ Réussi \\
Paiement Stripe & ✓ Réussi \\
Génération factures & ✓ Réussi \\
Dashboard Admin & ✓ Réussi \\
\hline
\end{tabular}
\end{center}

% ================================
% 6. INTERFACE UTILISATEUR
% ================================
\section{Interface utilisateur}

\subsection{Design et ergonomie}
\begin{itemize}
    \item Design moderne avec Bootstrap 5
    \item Utilisation de Font Awesome pour les icônes
    \item Animations CSS personnalisées
    \item Dégradés de couleurs pour les boutons et alertes
\end{itemize}

\subsection{Responsive design}
L'application est entièrement responsive et testée sur :
\begin{itemize}
    \item Desktop (1920×1080)
    \item Tablette (768×1024)
    \item Mobile (375×667)
\end{itemize}

\subsection{Captures d'écran}
% Insérer captures d'écran ici
\begin{figure}[h]
\centering
% \includegraphics[width=0.8\textwidth]{images/accueil.png}
\caption{Page d'accueil avec liste des créneaux}
\end{figure}

\begin{figure}[h]
\centering
% \includegraphics[width=0.8\textwidth]{images/panier.png}
\caption{Page du panier}
\end{figure}

\begin{figure}[h]
\centering
% \includegraphics[width=0.8\textwidth]{images/admin.png}
\caption{Dashboard administrateur}
\end{figure}

% [Ajouter d'autres captures]

% ================================
% 7. DIFFICULTÉS ET SOLUTIONS
% ================================
\section{Difficultés rencontrées et solutions}

\subsection{Intégration Stripe}
\textbf{Difficulté:} Configuration initiale de Stripe Elements et gestion asynchrone.\\
\textbf{Solution:} Utilisation de la dernière version de Stripe.js et gestion appropriée des promesses JavaScript.

\subsection{Gestion des places disponibles}
\textbf{Difficulté:} Problèmes de concurrence lors de réservations simultanées.\\
\textbf{Solution:} Validation côté serveur avec transactions et vérification des places avant confirmation.

% [Ajouter d'autres difficultés rencontrées]

% ================================
% 8. AMÉLIORATIONS FUTURES
% ================================
\section{Améliorations possibles}

\begin{enumerate}
    \item \textbf{Notifications:} Envoi d'emails de confirmation et rappels
    \item \textbf{Calendrier:} Vue calendrier pour les créneaux
    \item \textbf{Recherche avancée:} Filtres supplémentaires (prix, notes)
    \item \textbf{Système d'avis:} Permettre aux clients de noter les terrains
    \item \textbf{Gestion des équipes:} Créer et gérer des équipes
    \item \textbf{Statistiques avancées:} Graphiques interactifs avec Chart.js
    \item \textbf{Export de données:} Export Excel/PDF des rapports
    \item \textbf{API REST:} Créer une API pour applications mobiles
    \item \textbf{Internationalisation:} Support multilingue
    \item \textbf{Dark mode:} Thème sombre
\end{enumerate}

% ================================
% 9. CONCLUSION
% ================================
\section{Conclusion}

Ce projet m'a permis de développer une application web complète en utilisant les technologies ASP.NET Core MVC. J'ai pu mettre en pratique :

\begin{itemize}
    \item La conception d'une architecture MVC robuste
    \item L'implémentation d'un système d'authentification multi-rôles
    \item L'intégration d'un système de paiement tiers (Stripe)
    \item La gestion de la persistance avec Entity Framework Core
    \item La création d'interfaces utilisateur modernes et responsives
\end{itemize}

Les objectifs du travail pratique ont été atteints avec succès, et l'application est fonctionnelle et prête à être déployée.

% ================================
% ANNEXES
% ================================
\appendix
\section{Annexe A : Structure du projet}

\begin{lstlisting}
TP1/
├── Controllers/
│   ├── AccountController.cs
│   ├── AdminController.cs
│   ├── CartController.cs
│   ├── CheckoutController.cs
│   ├── HomeController.cs
│   └── ReservationsController.cs
├── Models/
│   ├── ViewModels/
│   ├── Utilisateur.cs
│   ├── Terrain.cs
│   ├── Creneau.cs
│   ├── PanierItem.cs
│   ├── Reservation.cs
│   ├── Paiement.cs
│   └── Facture.cs
├── Views/
├── Services/
├── Data/
└── wwwroot/
\end{lstlisting}

\section{Annexe B : Comptes de test}

\begin{center}
\begin{tabular}{|l|l|l|}
\hline
\textbf{Rôle} & \textbf{Email} & \textbf{Mot de passe} \\
\hline
Admin & admin@terrains.com & Admin123! \\
Fournisseur & fournisseur@terrains.com & Fournisseur123! \\
Client & client@terrains.com & Client123! \\
\hline
\end{tabular}
\end{center}

\section{Annexe C : Instructions d'installation}

Voir le fichier \texttt{README.md} pour les instructions complètes d'installation et de configuration.

% ================================
% RÉFÉRENCES
% ================================
\begin{thebibliography}{9}

\bibitem{aspnet}
Microsoft,
\textit{ASP.NET Core Documentation},
\url{https://docs.microsoft.com/aspnet/core/}

\bibitem{ef}
Microsoft,
\textit{Entity Framework Core Documentation},
\url{https://docs.microsoft.com/ef/core/}

\bibitem{stripe}
Stripe,
\textit{Stripe API Documentation},
\url{https://stripe.com/docs/api}

\bibitem{bootstrap}
Bootstrap Team,
\textit{Bootstrap 5 Documentation},
\url{https://getbootstrap.com/docs/5.0/}

\end{thebibliography}

\end{document}

